% ====================================================================
% discussion.tex (dissertation example)
% ====================================================================

\section{Discussion}
\label{sec:discussion}

\subsection{Agreement with theory}

The sine wave test agreed with \cref{eq:enobTheory} to within 0.03 bits at
every oversampling ratio, which is within what would be expected from a run
of 4096 samples. This supports both the theory and the model, and it shows
that for a signal which is itself moving across many codes the method works
as the textbooks describe.

The static test is the more informative of the two, because it removes the
help that a moving signal gives. Here the theory held only when there was
enough noise. With \qty{0.5}{LSB} rms or more, the static results matched the
sine wave results, so at that level the noise alone is sufficient to make the
errors independent and the nature of the signal no longer matters.

\subsection{Why a small amount of noise is not enough}

The result at \qty{0.25}{LSB} can be explained by considering an input that
sits a quarter of the way between two codes. With plenty of noise, the
readings are spread over several codes in proportions that average out to the
true value. With only a little, nearly all of the readings return the nearer
code, and the occasional reading of the further one does not occur often
enough to pull the average to the right place. The average then settles on a
value that is wrong by a fixed amount, and because the amount is fixed, more
averaging does not reduce it.

This is why the curve for \qty{0.25}{LSB} in \cref{fig:dcError} levels out
and does not continue downwards. The averaging is still removing the random
part of the error, but what remains is a systematic error that depends on
where the input sits within the step.

\subsection{The cost of too much noise}

The results also show that noise beyond what is needed is not free. At a
ratio of 256 the error with \qty{1.0}{LSB} of noise was almost twice that
with \qty{0.5}{LSB}, which is a loss of nearly one bit, and to recover it the
ratio would have to be raised by a factor of four. The middle term of
\cref{eq:enobTheory} predicts exactly this. The guidance to provide at
least one LSB of noise is therefore safe, but it is not optimal, and a
designer who can control the noise level would do better to aim for about
half an LSB.

\subsection{Practical implications}

For a 12-bit converter with a \qty{3.3}{\volt} reference, half an LSB is
\qty{0.4}{\milli\volt} rms. Many circuits will have this much noise without
any being added, in which case the method can be used as it is. Whether a
particular circuit does can be checked very simply, by holding the input
steady and looking at the spread of the raw codes. If nearly all of them are
the same code, there is too little noise and dither will have to be added.

The other practical cost is time. An output rate of \qty{100}{\hertz} with a
ratio of 256 needs the converter to run at \qty{25.6}{\kilo\hertz}, which is
well within the reach of a typical microcontroller, but the same ratio at an
output rate of \qty{10}{\kilo\hertz} would need \qty{2.56}{\mega\hertz},
which is not.

\subsection{Limitations}

The model is of an ideal converter, and three of its simplifications affect
how far the results can be carried over to a real one.

The first is that the steps of the model are all exactly the same size. A
real converter has steps of slightly different sizes, described by its
differential and integral non-linearity, and averaging does nothing to
correct these, as they are the same on every reading. The 15 bits reached
here should be read as an upper limit, with the real figure set by the
linearity of the device in question.

The second is that the noise in the model is Gaussian and unrelated from one
sample to the next. Noise picked up from a switching power supply or from the
mains is neither, and noise that repeats in step with the sampling would not
average away at all.

The third is that a plain average is a poor low-pass filter. It was used
because it is what most firmware does, but a signal with content close to
half the output rate would be attenuated by it, and a better decimation
filter would be needed in that case~\cite{oppenheim_schafer_2010}.
